\clearpage
\VerdierPhysicalPageAnchor{0085}%
\VerdierSemanticAnchor{ch1-3-1-4-homotopy-definition-continuation}
deux morphismes de $\operatorname{tr}_1(C)$. Une \emph{homotopie} reliant
$(g,h,a)$ à $(g',h',a')$ est un triplet $(b,b',\alpha)$, où
$b\colon g\rightsquigarrow g'$ et $b'\colon h\rightsquigarrow h'$ sont des
homotopies et où

\VerdierSemanticAnchor{ch1-3-1-4-second-order-condition}
\begin{equation*}
  \alpha\colon f'\star b+a'\rightsquigarrow b'\star f+a
\end{equation*}
est une homotopie du deuxième ordre (2.5.12).

\VerdierSemanticAnchor{ch1-3-1-4-homotopic-morphisms}
Deux morphismes de $\operatorname{tr}_1(C)$ seront dits \emph{homotopes} s'il
existe une homotopie les reliant.
\VerdierSemanticAnchor{ch1-3-1-4-equivalence-relation}
La relation : ``$(g,h,a)$ et $(g',h',a')$ sont homotopes'' est une relation
d'équivalence compatible avec la somme directe et la composition des
morphismes.
\VerdierSemanticAnchor{ch1-3-1-4-tr1-homotopy-category}
On peut donc introduire la catégorie $\operatorname{Tr}_1(C)$ dont les objets
sont les objets de $\operatorname{tr}_1(C)$ et les morphismes les morphismes de
$\operatorname{tr}_1(C)$ à homotopie près.

\VerdierSemanticAnchor{ch1-3-1-5}
\noindent\textbf{Proposition 3.1.5.}\quad
\emph{Soient
$D,D'\colon (X\xrightarrow{f}Y)\rightrightarrows
(X'\xrightarrow{f'}Y')$ deux morphismes de $\operatorname{tr}_1(C)$. Si les
morphismes $D$ et $D'$ sont homotopes dans $\operatorname{tr}_1(C)$, alors les
morphismes $c(D)$ et $c(D')$ sont homotopes dans $\operatorname{comp}(C)$.}

\VerdierSemanticAnchor{ch1-3-1-5-proof-setup}
Soient $D=(g,h,a)$ et $D'=(g',h',a')$. Supposons que $D$ et $D'$ soient
homotopes dans $\operatorname{tr}_1(C)$. Il existe donc des homotopies $b$ et
$b'$ telles que :

\VerdierSemanticAnchor{ch1-3-1-5-eq1}
\begin{equation}
  \begin{aligned}
    g'-g&=bd+db,\\
    h'-h&=b'd+db'.
  \end{aligned}
  \tag{3.1.5.1}
\end{equation}

(Pour alléger la notation, nous avons omis dans ce calcul les exposants
indiquant le degré des composants.)
\VerdierSemanticAnchor{ch1-3-1-5-second-order-homotopy}
Il existe de plus une homotopie du deuxième ordre $\alpha$ telle que :
\begin{equation*}
  \alpha d-d\alpha=b'f+a-(f'b+a').
\end{equation*}

\VerdierSemanticAnchor{ch1-3-1-5-cone-homotopy}
Un calcul immédiat montre alors que
\begin{equation*}
  \begin{pmatrix}
    b' & \alpha\\
    0  & -b
  \end{pmatrix}
\end{equation*}
est une homotopie qui relie $c(D)$ à $c(D')$.

\VerdierSemanticAnchor{ch1-3-1-6}
\noindent\textbf{Corollaire 3.1.6.}\quad
\emph{Les foncteurs :}
\begin{equation*}
  c\colon\operatorname{tr}_1(C)\longrightarrow\operatorname{comp}(C)
  \qquad\text{et}\qquad
  s\colon\operatorname{comp}(C)\longrightarrow\operatorname{tr}_1(C)
\end{equation*}
\emph{sont compatibles avec les homotopies de $\operatorname{tr}_1(C)$ et de
$\operatorname{comp}(C)$ et définissent, en passant au quotient, des foncteurs
encore notés}

\VerdierSemanticAnchor{ch1-3-1-6-descended-functors}
\begin{equation*}
  c\colon\operatorname{Tr}_1(C)\longrightarrow K(C)
  \qquad\text{et}\qquad
  s\colon K(C)\longrightarrow\operatorname{Tr}_1(C).
\end{equation*}
\VerdierSemanticAnchor{ch1-3-1-6-adjunction}
\emph{Le foncteur $c\colon\operatorname{Tr}_1(C)\longrightarrow K(C)$ est
adjoint à gauche du foncteur
$s\colon K(C)\longrightarrow\operatorname{Tr}_1(C)$.}

\VerdierSemanticAnchor{ch1-3-1-6-proof}
Ceci résulte formellement de (3.1.5) et de (3.1.3).

\VerdierSemanticAnchor{ch1-3-1-7}
\noindent\textbf{3.1.7.}\quad
\VerdierSemanticAnchor{ch1-3-1-7-s-prime}
Désignons par
$s'\colon\operatorname{comp}(C)\longrightarrow\operatorname{tr}_1(C)$ le
foncteur $X\longmapsto(X\longrightarrow 0)$.
\VerdierSemanticAnchor{ch1-3-1-7-comp-translation}
Le foncteur
$cs'\colon\operatorname{comp}(C)\longrightarrow\operatorname{comp}(C)$ est
égal au foncteur de translation de $\operatorname{comp}(C)$.
\VerdierSemanticAnchor{ch1-3-1-7-descended-s-prime}
Le foncteur $s'$ transforme deux morphismes homotopes de
$\operatorname{comp}(C)$ en morphismes homotopes de $\operatorname{tr}_1(C)$ et
par suite définit un foncteur que nous noterons encore
$s'\colon K(C)\longrightarrow\operatorname{Tr}_1(C)$.
\VerdierSemanticAnchor{ch1-3-1-7-k-translation}
Le foncteur $cs'\colon K(C)\longrightarrow K(C)$ est égal au foncteur de
translation de $K(C)$.
